\subsection{FORM ASSOCIATED TO A COXETER GROUP}

Let \(S\) be a set and let \(M = (m(s, s')_{s, s' \in S})\) be a Coxeter matrix (Chap. IV, §1, no. 9) of type \(S\) . Recall that this means:

\begin{enumerate}
\def\labelenumi{(\arabic{enumi})}
\item
  The elements of \(M\) are integers or \(+\infty\) .
\item
  M is symmetric.
\item
  \(m(s, s) = 1\) for all s.
\item
  \(m(s, s') \geqslant 2\) for \(s \neq s'\) .
\end{enumerate}

Let \(\mathbf{E} = \mathbf{R}^{(\mathrm{S})}\) , let \((e_s)_{s \in S}\) be the canonical basis of \(\mathbf{E}\) , and let \(B_M\) be the bilinear form on \(\mathbf{E}\) such that

\[
\mathrm{B} _ {\mathrm{M}} (e _ {s}, e _ {s ^ {\prime}}) = - \cos \frac {\pi}{m (s , s ^ {\prime})}.
\]

The form \(B_{M}\) is symmetric. It is called the associated form of the matrix M. We have

\[
\mathrm{B} _ {\mathrm{M}} \left(e _ {s}, e _ {s}\right) = 1 \text {and} \mathrm{B} _ {\mathrm{M}} \left(e _ {s}, e _ {s ^ {\prime}}\right) \leqslant 0 \text {if} s \neq s ^ {\prime}.
\]

Let \(s \in S\) and let \(f_s\) be the linear form \(x \mapsto 2B_M(e_s, x)\) . We denote by \(\sigma_s\) the pseudo-reflection defined by the pair \((e_s, f_s)\) (cf.~§2, no. 1); since \(\langle e_s, f_s \rangle = 2\) , it is a reflection (§2, no. 2). We have

\[
\sigma_ {s} (x) = x - 2 \mathrm{B} _ {\mathrm{M}} (e _ {s}, x). e _ {s}
\]

and in particular

\[
\sigma_ {s} (e _ {s ^ {\prime}}) = e _ {s ^ {\prime}} + 2 \cos \frac {\pi}{m (s , s ^ {\prime})}. e _ {s}.
\]

Since \(e_{s}\) is not isotropic for \(B_{M}\) , the space E is the direct sum of the line \(Re_{s}\) and the hyperplane \(H_{s}\) orthogonal to \(e_{s}\) . Since \(\sigma_{s}\) is equal to -1 on \(Re_{s}\) and to 1 on \(H_{s}\) , it follows that \(\sigma_{s}\) preserves the form \(B_{M}\) . When S is finite and \(B_{M}\) is non-degenerate (a case to which we shall return in no. 8), it follows that \(\sigma_{s}\) is an orthogonal reflection ( \(\S2\) , no. 3).

